\chapter{系统测试}

\section{测试环境与工具}

系统测试的目标在于验证本文所实现的企业供应链协同管理系统是否能够在真实业务语义下正确完成认证授权、订单流转、库存变化、生产协同、采购协同、质检回流、质量追溯和实时通知等功能。由于系统属于典型的前后端分离项目，测试不应仅关注后端服务是否返回正确结果，还需要兼顾前端接口联调和多角色流程闭环验证。

结合当前项目实现，系统测试环境如表~\ref{tab:test-env-real} 所示。

\begin{table}[htbp]
\centering
\footnotesize
\caption{系统测试环境与工具配置}
\label{tab:test-env-real}
\begin{tabular}{|c|c|p{6.2cm}|}
\hline
类别 & 配置项 & 说明 \\\hline
操作系统 & Windows & 本地后端、前端与 LaTeX 编译环境。 \\\hline
JDK & Java 17 & Spring Boot 后端运行及单元测试环境。 \\\hline
Maven & Maven 3.x & 后端依赖管理、构建与测试执行。 \\\hline
Node.js & Node.js 运行环境 & 执行前端接口联通、实时消息和端到端脚本。 \\\hline
数据库 & MySQL & 存储系统业务数据，支撑接口联调与流程验证。 \\\hline
前端环境 & Vue 3 + Vite & 提供页面运行与交互验证环境。 \\\hline
浏览器 & Chrome / Edge & 验证角色页面、订单流程和消息接收效果。 \\\hline
后端测试框架 & JUnit 5、Mockito & 验证服务逻辑、控制器行为与状态流转规则。 \\\hline
前端联调脚本 & \texttt{test:api}、\texttt{test:realtime}、\texttt{test:e2e}、\texttt{test:workflow} & 验证接口可达性、实时消息可达性和关键业务联调流程。 \\\hline
\end{tabular}
\end{table}

根据项目当前目录结构，测试主要由两部分组成：一部分位于后端 \texttt{src/test/java} 下，以 JUnit 5 和 Mockito 为基础，对工作流服务、周计划服务、质检服务以及部分控制器进行验证；另一部分位于前端 \texttt{frontend/scripts} 下，以 Node.js 脚本形式对 API、实时连接和端到端流程进行联调测试。这种组合式测试方案与系统的前后端分离结构相匹配，能够较全面地反映项目实现质量。

\section{单元测试}

\subsection{后端单元测试覆盖范围}

当前项目已经围绕认证、订单、采购、质检和周计划等核心模块编写了后端测试类，主要包括 \texttt{AuthControllerTest}、\texttt{CustomerControllerTest}、\texttt{SalesOrderControllerTest}、\texttt{ProductionControllerTest}、\texttt{ProcurementControllerTest}、\texttt{QualityControllerTest}、\texttt{OrderWorkflowServiceTest}、\texttt{ProcurementWorkflowServiceTest}、\texttt{QualityServiceTest} 和 \texttt{WeeklyPlanningServiceTest} 等。这些测试覆盖了核心业务对象的状态变化和关键路径校验，是验证系统正确性的重要依据。

与仅测试“接口是否返回 200”不同，当前后端测试更关注以下问题：第一，状态迁移是否符合业务规则；第二，异常输入是否能够被拒绝；第三，关键业务动作是否触发消息通知；第四，库存、批次和计划等派生对象是否被正确创建或更新。由此可见，测试目标并不仅限于程序语法正确，而是深入到业务一致性层面。

\subsection{订单工作流测试}

订单工作流测试主要对应 \texttt{OrderWorkflowServiceTest}。该测试覆盖了仓库核查库存充足、库存不足生成生产计划、生产完工转待质检、质检合格后仓库确认生产入库以及发货通知顾客等关键场景。例如，在“库存充足”场景下，测试会构造订单与库存数据，验证系统是否正确增加预留库存、是否将订单状态推进为“已接单”，以及是否向销售主题发送库存核验通过通知；在“库存不足”场景下，测试验证系统是否生成生产计划、订单是否转入“生产中”，以及是否向生产主题发送通知。

这些测试说明，订单工作流已不再是简单的状态字段更新，而是伴随着库存预留、计划生成、消息推送等一系列联动操作。通过对这些联动效果进行断言，可以较好地保证订单协同链路在后续迭代中保持稳定。

\subsection{采购协同测试}

采购模块的核心测试集中在 \texttt{ProcurementWorkflowServiceTest}。测试重点包括：采购单创建后是否进入“待供应商确认”状态，是否自动生成采购单号并向供应商主题发送通知；供应商发货后是否正确通知采购管理员；仓库收货后是否完成库存累加与状态变更；采购单创建时是否校验所选供应商与原材料档案中的供应商信息一致。

其中，“供应商与原材料不匹配时拒绝创建采购单”的测试尤为重要。该测试构造了优选供应商为“供应商B”的原材料，而采购管理员却选择“供应商A”，系统应返回 400 状态并提示原材料与所选供应商不匹配。这一结果说明采购模块在实现过程中真正贯彻了主数据约束，而不是在表单层面做形式性限制。

\subsection{质检服务测试}

\texttt{QualityServiceTest} 主要验证批次质检后的不同分支处理逻辑。测试内容包括：当质检结果为不合格且填写原因时，系统是否向对应生产管理员主题发送返工通知、是否写入质检记录并标记通知已发送；当质检结果为合格时，系统是否仅保存记录而不通知生产管理员；当质检结果为不合格但未填写原因时，系统是否拒绝该请求。通过这些测试，可以确保质检模块既满足业务严谨性，也能对异常路径进行有效约束。

此外，质检服务还包含对“查看当前生产管理员的不合格批次告警”功能的测试。这说明系统不仅支持结果写入，还支持基于责任人反查异常质量问题，增强了生产—质检之间的闭环性。

\subsection{周计划服务测试}

\texttt{WeeklyPlanningServiceTest} 对生产周计划和采购周计划生成逻辑进行了针对性验证。测试并非只断言计划对象非空，而是进一步验证在给定上周完工量、订单需求、库存数量、BOM 关系和在途采购量后，系统是否能按照预定公式计算出正确的建议数量。例如，生产计划测试验证了“上周完工量 10、当前需求 8、库存 3”时的基线与建议生产量；采购计划测试验证了 BOM 需求、上周采购量、库存和在途量共同作用下的建议采购量是否正确。

该类测试表明，周计划模块虽然算法复杂度不高，但其规则表达具有明确的可验证性。通过单元测试固化这些规则，可以显著降低后续功能调整时引入计算错误的风险。

\subsection{单元测试汇总}

为更直观地说明后端单元测试覆盖内容，本文将代表性测试整理如表~\ref{tab:unit-test-real} 所示。

\begin{table}[htbp]
\centering
\footnotesize
\caption{代表性后端单元测试内容}
\label{tab:unit-test-real}
\begin{tabular}{|c|p{4.2cm}|p{5.8cm}|}
\hline
测试类 & 主要验证内容 & 预期结果 \\\hline
OrderWorkflowServiceTest & 仓库核查、生产计划生成、发货通知、生产入库确认 & 状态迁移正确，库存和消息联动正确。 \\\hline
ProcurementWorkflowServiceTest & 采购单创建、供应商发货、仓库收货、供应商匹配校验 & 采购状态正确推进，不匹配供应商被拒绝。 \\\hline
QualityServiceTest & 合格/不合格质检、返工通知、不合格原因必填 & 不同分支处理逻辑符合业务要求。 \\\hline
WeeklyPlanningServiceTest & 生产周计划、采购周计划公式计算与通知 & 建议数量与规则一致，计划生成后触发通知。 \\\hline
SalesOrderControllerTest 等控制器测试 & 接口入参处理、权限与返回结构 & 控制器能正确委派服务并返回规范结果。 \\\hline
\end{tabular}
\end{table}

综合来看，现有单元测试已经覆盖了系统的多条关键业务主线，为论文后续功能测试与结果分析提供了可信支撑。

\section{功能测试}

\subsection{测试思路}

功能测试以“多角色协同闭环是否成立”为中心展开，不再局限于某一个接口返回是否正确，而是观察一个业务对象在多个角色处理后能否按照预定流程持续推进。由于本文系统的最大特点在于角色之间的衔接关系，因此功能测试的重点应放在订单、采购和质检三条主流程上。

具体测试思路如下：首先验证登录与权限隔离是否有效，确保不同角色只能看到自身功能范围；其次验证顾客下单、销售审核、仓库核查、生产完工、质检判定、仓库入库与发货这一完整订单链路；再次验证采购管理员创建采购单、供应商接单发货、采购管理员通知仓库、仓库收货入库这一采购协同链路；最后验证质量追溯、库存预警和周计划建议是否能够输出合理结果。

\subsection{前端联调脚本与接口验证}

项目当前前端提供了多项联调脚本，分别用于验证不同层面的系统可用性：\texttt{npm run test:api} 用于验证后端接口可达性；\texttt{npm run test:realtime} 用于验证 WebSocket 实时消息端点可达性；\texttt{npm run test:e2e} 用于验证注册、登录、创建订单及接收订单主题消息的基本闭环；\texttt{npm run test:workflow} 用于验证更完整的客户下单—订单处理工作流。尽管这些脚本本质上属于工程验证工具，但它们为论文中的功能测试提供了较好的自动化基础。

在当前项目产物中，还保留了周计划相关的联调结果文件 \texttt{weekly-plan-smoke-summary.json}。该文件显示，在一次以 2026 年 4 月 16 日为基准的验证中，生产管理员与采购管理员测试账号均能注册和登录成功，生产周计划列表接口与当前计划接口均返回 200 状态，采购周计划列表接口与当前计划接口同样返回 200 状态，其中当前生产周计划明细数为 1，当前采购周计划明细数为 5。这一结果说明，周计划接口在当前实现中已能完成基本联调并返回结构化结果。

\subsection{主要功能场景测试}

结合系统功能，本文将代表性功能测试场景整理如表~\ref{tab:function-test-real} 所示。

\begin{table}[htbp]
\centering
\footnotesize
\caption{主要功能测试场景与结果}
\label{tab:function-test-real}
\begin{tabular}{|c|p{5.4cm}|c|}
\hline
测试场景 & 验证要点 & 结果 \\\hline
登录与角色隔离 & 不同角色登录后菜单与接口权限范围正确，越权访问被拒绝 & 通过 \\\hline
顾客下单与销售审核 & 顾客创建订单后进入待审核，销售管理员可接受或拒绝 & 通过 \\\hline
仓库核查与发货 & 库存充足时完成预留和发货，库存不足时生成生产计划 & 通过 \\\hline
生产完工与质检回流 & 生产完工后生成待检批次，质检不合格时退回生产 & 通过 \\\hline
生产入库确认 & 仅质检合格批次允许仓库执行生产入库确认 & 通过 \\\hline
采购与供应商协同 & 采购单自动生成编号，供应商可接单/拒单/发货，仓库可收货入库 & 通过 \\\hline
原材料档案约束 & 采购时供应商必须与原材料档案中维护的信息匹配 & 通过 \\\hline
质量追溯 & 顾客可根据订单号查看对应质检结果 & 通过 \\\hline
库存预警与周计划 & 系统可识别成品与原材料缺口并给出计划建议 & 通过 \\\hline
实时通知 & 订单、采购、质检等关键事件可推送给对应角色主题 & 通过 \\\hline
\end{tabular}
\end{table}

从功能测试结果来看，系统已经能够支撑论文中提出的大部分核心业务场景，尤其是在多角色闭环协同和状态同步方面表现较为完整。系统中“一个角色处理完成后触发下一个角色待办”的机制已初步成型，这也是本文系统区别于静态信息录入系统的关键特征。

\section{测试结果分析}

综合单元测试和功能测试结果，可以从以下几个方面分析系统表现。

首先，从功能完整性角度看，系统已基本实现顾客下单、销售审核、仓库核查、生产协同、采购协同、质检处理、质量追溯、库存预警和周计划生成等主要模块，且这些模块并非彼此孤立，而是通过订单号、批次号、采购单号和消息主题形成了较为清晰的业务链路。特别是在订单—仓储—生产—质检闭环和采购—供应商—仓库闭环两条主线上，系统能够根据库存状态和质检结论自动调整后续流程，体现出较好的业务一致性。

其次，从数据一致性角度看，测试结果表明系统在关键状态更新时通常伴随着关联对象同步变化。例如，仓库发货不仅会将订单状态更新为已发货，还会扣减库存并生成出库流水；采购收货不仅更新采购单状态，还会增加原材料库存并写入入库记录；质检不合格不仅写入质检记录，还会将订单退回生产状态。这说明项目在设计时已经将“状态变化”和“数据留痕”一并考虑，能够较好地支持审计与追溯。

再次，从实时协同能力角度看，WebSocket/STOMP 消息机制已在多个关键节点得到应用。测试中能够验证订单主题、采购主题、质检主题和周计划主题的消息广播行为，这表明系统已经具备事件驱动式协同的基本能力。相较于传统只依赖轮询刷新的方式，这种实现更符合企业多岗位协作的实际需求。

当然，测试结果也暴露出系统仍有进一步优化空间。其一，当前周计划生成虽然具备可解释性，但仍以规则计算为主，对复杂需求波动、季节因素和多约束资源排产的适应能力有限；其二，前端联调脚本已覆盖接口和实时通信层面，但针对高并发、多用户同时在线的性能测试仍有待补充；其三，当前系统更偏向 Web 端，若未来面向移动巡检、移动仓储作业或移动审批场景扩展，还需补充针对移动端体验和离线容错能力的测试。

总体来看，本文系统已能够满足毕业设计阶段对于供应链协同管理系统的主要功能要求，并在多角色业务闭环、实时消息同步和质量追溯等方面体现出较好的工程实现效果。测试结果证明，该系统不仅能够完成基础业务数据管理，还能够围绕订单履约和物料保障建立较为可靠的协同机制，为后续功能优化和实际应用扩展奠定了基础。

